% \iffalse meta-comment
% Copyright (C) 2026 Andreas Willibald Weber
%
% This work may be distributed and/or modified under the conditions of the
% LaTeX Project Public License, either version 1.3c of this license or (at your
% option) any later version. The latest version of the license is available at
% https://www.latex-project.org/lppl.txt.
%
% This work has the LPPL maintenance status `maintained'. The Current
% Maintainer of this work is Andreas Willibald Weber.
%
% This file is part of the `kanban' package; the complete list of files
% constituting the work appears in the main source file kanban.dtx.
% \fi
%
%    \begin{macrocode}
%<*package>
%    \end{macrocode}
%
% \subsection{Persistent configuration}
%
% Built-in palettes are registered through the public mechanism. Mode, date
% format, start date, reference date, column list, palette, semantic labels,
% and gap defaults then enter the dedicated |kanban/setup| LaTeX3 key family.
% Package options are accumulated and passed through the same transaction so
% package-level and document-level setup have identical validation and scope.
% The functions in this subsection distinguish persistent assignment from
% board-local assignment consistently: a package option or |\kbset| changes the
% source for later snapshots, whereas an environment key changes only the
% snapshot already being constructed.
%
% \paragraph{Built-in palette registration.}
% The package's default, muted, and monochrome palettes should obey exactly the
% same validation and registry rules as user-defined palettes. They are
% therefore registered through the canonical public command rather than by
% writing directly to internal property lists. The default palette is selected
% immediately after registration so a document with no configuration still
% produces the documented colors.
%
% These are ordinary calls to |\kanbandefinepalette|, whose two mandatory
% arguments provide the registry name and comma list. The percent sign after
% the opening brace suppresses an otherwise unwanted source-line space in the
% argument. |\kanban@selectpalette| then invokes the validated persistent
% selector described in the preceding subsection. Reusing public behavior in
% this way makes built-in and user palettes follow one code path.
%    \begin{macrocode}
\kanbandefinepalette{default}{%
  kanbanSlate,kanbanBlue,kanbanTeal,kanbanSage,
  kanbanAmber,kanbanClay,kanbanRose,kanbanPlum}
\kanbandefinepalette{muted}{kanbanMutedBlue,kanbanMutedGreen,
  kanbanMutedPurple,kanbanMutedBrown}
\kanbandefinepalette{monochrome}{black!85,black!75,black!65,black!55}
\kanban@selectpalette{default}
%    \end{macrocode}
%
% \paragraph{Mode selection.}
% Visual mode is a constrained textual choice. The persistent selector records
% a global |light| or |dark| default; the board selector records the same choice
% locally. Unknown values do not reset the previous selection.
%
% |\newcommand| defines a checked LaTeX2e command and |[1]| declares one
% argument. Both selectors first call the shared lowercase choice normalizer.
% |\ifdefstring| from |etoolbox| compares the replacement text of a macro with
% a literal string and accepts true and false branches. The global selector
% uses |\gdef|, while the board selector uses |\def| so its result disappears
% with the environment group. The final false branch emits a recoverable
% |\PackageWarning|.
%    \begin{macrocode}
\newcommand\kanban@selectmode[1]{%
  \kanban@normalizechoice{#1}%
  \ifdefstring{\kb@choice}{light}{\gdef\kanban@mode{light}}{%
    \ifdefstring{\kb@choice}{dark}{\gdef\kanban@mode{dark}}{%
      \PackageWarning{kanban}
        {unknown mode '#1'; keeping the current mode}}}}
\newcommand\kanban@selectboardmode[1]{%
  \kanban@normalizechoice{#1}%
  \ifdefstring{\kb@choice}{light}{\def\kanban@boardmode{light}}{%
    \ifdefstring{\kb@choice}{dark}{\def\kanban@boardmode{dark}}{%
      \PackageWarning{kanban}
        {unknown mode '#1'; keeping the current mode}}}}
%    \end{macrocode}
%
% \paragraph{Date-format canonicalization and selection.}
% Five canonical format names are supported, together with a small set of
% readable aliases. Canonicalization is separated from assignment so the same
% mapping serves both persistent and board-local selection. A recognized alias
% writes one of |us|, |english|, |german|, |iso|, or |ymd|; an unknown choice
% leaves the result empty and is diagnosed by the appropriate selector.
%
% |\def\kb@datechoice{}| initializes the result sentinel. Repeated
% |\ifdefstring| tests implement the alias table, and a successful test replaces
% the sentinel with its canonical value. |\ifdefempty| then distinguishes
% failure from success. |\global\let| assigns the persistent format macro the
% current meaning of the canonical result without retaining an intermediate
% expansion dependency; the local |\let| does the same for a board snapshot.
% |\MessageBreak| improves the line wrapping of the longer board warning.
%    \begin{macrocode}
\newcommand\kanban@canonicaldateformat[1]{%
  \kanban@normalizechoice{#1}%
  \def\kb@datechoice{}%
  \ifdefstring{\kb@choice}{us}{\def\kb@datechoice{us}}{}%
  \ifdefstring{\kb@choice}{american}{\def\kb@datechoice{us}}{}%
  \ifdefstring{\kb@choice}{english}{\def\kb@datechoice{english}}{}%
  \ifdefstring{\kb@choice}{uk}{\def\kb@datechoice{english}}{}%
  \ifdefstring{\kb@choice}{british}{\def\kb@datechoice{english}}{}%
  \ifdefstring{\kb@choice}{german}{\def\kb@datechoice{german}}{}%
  \ifdefstring{\kb@choice}{de}{\def\kb@datechoice{german}}{}%
  \ifdefstring{\kb@choice}{deutsch}{\def\kb@datechoice{german}}{}%
  \ifdefstring{\kb@choice}{iso}{\def\kb@datechoice{iso}}{}%
  \ifdefstring{\kb@choice}{ymd}{\def\kb@datechoice{ymd}}{}%
  \ifdefstring{\kb@choice}{year-first}{\def\kb@datechoice{ymd}}{}}
\newcommand\kanban@selectdateformat[1]{%
  \kanban@canonicaldateformat{#1}%
  \ifdefempty\kb@datechoice
    {\PackageWarning{kanban}
      {unknown date format '#1'; keeping the current date format}}
    {\global\let\kanban@dateformat\kb@datechoice}}
\newcommand\kanban@selectboarddateformat[1]{%
  \kanban@canonicaldateformat{#1}%
  \ifdefempty\kb@datechoice
    {\PackageWarning{kanban}
      {unknown board date format '#1';
        \MessageBreak keeping the current date format}}
    {\let\kanban@boarddateformat\kb@datechoice}}
%    \end{macrocode}
%
% \paragraph{Validated column-gap assignment.}
% A column gap must be a non-negative dimension. The persistent and board-local
% setters share the safe grammar established by the deadline subsystem, assign
% the accepted input to a candidate register, and compare the evaluated length
% with zero before committing it to the appropriate scope.
%
% |\kanban@validatedimensionTF| supplies success and failure branches.
% |\setlength| performs the actual LaTeX length assignment only after syntax
% validation. The primitive |\ifdim| compares dimensions; its true branch for a
% negative candidate issues the shared warning. A successful persistent update
% uses |\global| before the register assignment, while the board update remains
% local. The closing |\fi| terminates the primitive conditional.
%    \begin{macrocode}
\newcommand\kanban@setcolumnsep[1]{%
  \kanban@validatedimensionTF{#1}{%
    \setlength\kanban@candidatecolumnsep{#1}%
    \ifdim\kanban@candidatecolumnsep<0pt
      \kanban@dimensionwarning{columnsep}{#1}%
    \else
      \global\kanban@columnsep=\kanban@candidatecolumnsep
    \fi
  }{%
    \kanban@dimensionwarning{columnsep}{#1}}}
\newcommand\kanban@setboardcolumnsep[1]{%
  \kanban@validatedimensionTF{#1}{%
    \setlength\kanban@candidatecolumnsep{#1}%
    \ifdim\kanban@candidatecolumnsep<0pt
      \kanban@dimensionwarning{board columnsep}{#1}%
    \else
      \kanban@boardcolumnsep=\kanban@candidatecolumnsep
    \fi
  }{%
    \kanban@dimensionwarning{board columnsep}{#1}}}
%    \end{macrocode}
%
% \paragraph{Materialization of the visual mode.}
% The mode name is semantic state; rendering requires concrete text, surface,
% border, and default heading colors. The application function translates the
% active board mode into those values after all board options have been
% processed. Metadata chip colors are intentionally absent from this function
% because their documented light-theme styling remains stable in both modes.
%
% |\ifdefstring| selects the dark branch when the board-mode macro expands to
% |dark| and otherwise uses the light definitions. Each |\def| is local to the
% board environment. The stored values are named |xcolor| colors or mix
% expressions; later |tcolorbox| definitions expand them when drawing a card.
%    \begin{macrocode}
\newcommand\kanban@applymode{%
  \ifdefstring{\kanban@boardmode}{dark}{%
    \def\kanban@textcolor{white!88!black}%
    \def\kanban@cardcolor{kanbanInk}%
    \def\kanban@bordercolor{white!18!black}%
    \def\kanban@headingtext{white}%
  }{%
    \def\kanban@textcolor{kanbanInk}%
    \def\kanban@cardcolor{white}%
    \def\kanban@bordercolor{black!15}%
    \def\kanban@headingtext{white}%
  }}
%    \end{macrocode}
%
% \paragraph{Persistent and board-local semantic labels.}
% Priority and completion labels are customizable text, but an empty label
% would create an invisible chip with misleading spacing. The two setters
% therefore reject blank input and dynamically address either the persistent
% or board-local label family.
%
% |\ifblank| from |etoolbox| tests the argument without requiring it to be a
% predefined macro. |\csgdef| globally defines a control sequence whose name is
% assembled from character data, such as |kanban@highlabel|. |\csdef| performs
% the analogous local definition for names such as
% |kanban@boardhighlabel|. Dynamic |cs| forms keep four nearly identical label
% setters in one implementation while the key family supplies the validated
% label identifier.
%    \begin{macrocode}
\newcommand\kanban@setlabel[2]{%
  \ifblank{#2}
    {\PackageWarning{kanban}
      {label '#1' must not be empty; keeping the current value}}
    {\csgdef{kanban@#1label}{#2}}}
\newcommand\kanban@setboardlabel[2]{%
  \ifblank{#2}
    {\PackageWarning{kanban}
      {board label '#1' must not be empty; keeping the current value}}
    {\csdef{kanban@board#1label}{#2}}}
%    \end{macrocode}
%
% \paragraph{Nested label key families.}
% Label customization is expressed as a subordinate key list, for example
% |labels={high=URGENT,done=FINISHED}|. Separate persistent and board-local
% families direct identical key names to setters with different assignment
% scope. Each family also owns its unknown-key diagnostic, which identifies
% whether the error occurred in global setup or in a board override.
%
% |\ExplSyntaxOn| is re-enabled for the LaTeX3 key definitions.
% |\keys_define:nn| creates or extends a named family. A property written
% |.code:n| stores executable code receiving the key value as |#1|. The
% reserved |unknown| handler receives the misspelled key name through
% |\l_keys_key_str| and issues the named diagnostic with
% |\msg_warning:nnnn|. Defining nested families independently also prevents a
% label name such as |mode| from being confused with a top-level setup key.
%    \begin{macrocode}
\ExplSyntaxOn
\keys_define:nn {kanban/setup-labels}
  {
    high .code:n = {\kanban@setlabel{high}{#1}},
    medium .code:n = {\kanban@setlabel{medium}{#1}},
    low .code:n = {\kanban@setlabel{low}{#1}},
    done .code:n = {\kanban@setlabel{done}{#1}},
    unknown .code:n =
      {\msg_warning:nnnn {kanban}{unknown-key}
        {label}{\l_keys_key_str}}
  }
\keys_define:nn {kanban/board-labels}
  {
    high .code:n = {\kanban@setboardlabel{high}{#1}},
    medium .code:n = {\kanban@setboardlabel{medium}{#1}},
    low .code:n = {\kanban@setboardlabel{low}{#1}},
    done .code:n = {\kanban@setboardlabel{done}{#1}},
    unknown .code:n =
      {\msg_warning:nnnn {kanban}{unknown-key}
        {board~label}{\l_keys_key_str}}
  }
%    \end{macrocode}
%
% \paragraph{The persistent setup key family.}
% The main setup family is the single semantic entry point for package options
% and later calls to |\kanbansetup|. Immediate keys delegate to their validated
% selectors. Date values are staged rather than applied immediately, because a
% |date-format| key appearing later in the same comma list must take effect
% before an ambiguous start or reference date is interpreted. Alias keys keep
% legacy or concise spelling without duplicating implementation.
%
% In |l3keys|, |.code:n| executes an arbitrary handler with the supplied value.
% |.meta:n| rewrites one key as another key assignment: |dateformat| becomes
% |date-format| and |today| becomes |reference-date|. The date handlers use
% |\tl_set:Nn| to retain the literal value and |\bool_set_true:N| to record that
% the staged value must later be committed. The |labels| key calls
% |\keys_set:nn| recursively on the subordinate family. The family-level
% |unknown| handler supplies a uniform non-fatal warning for unsupported setup
% names.
%    \begin{macrocode}
\keys_define:nn {kanban/setup}
  {
    columns .code:n = {\kanban@selectcolumns{#1}},
    palette .code:n = {\kanban@selectpalette{#1}},
    date-format .code:n = {\kanban@selectdateformat{#1}},
    dateformat .meta:n = {date-format={#1}},
    startdate .code:n =
      {
        \tl_set:Nn \l__kanban_pending_startdate_tl {#1}
        \bool_set_true:N \l__kanban_pending_startdate_bool
      },
    reference-date .code:n =
      {
        \tl_set:Nn \l__kanban_pending_reference_date_tl {#1}
        \bool_set_true:N \l__kanban_pending_reference_date_bool
      },
    today .meta:n = {reference-date={#1}},
    mode .code:n = {\kanban@selectmode{#1}},
    columnsep .code:n = {\kanban@setcolumnsep{#1}},
    labels .code:n = {\keys_set:nn {kanban/setup-labels}{#1}},
    unknown .code:n =
      {\msg_warning:nnnn {kanban}{unknown-key}
        {setup}{\l_keys_key_str}}
  }
%    \end{macrocode}
%
% \paragraph{A deferred setup transaction.}
% Processing a setup list is a small transaction. It clears both staging flags,
% lets |l3keys| process the complete list in source order, and only then applies
% any staged dates under the final selected format. Consequently
% |startdate=07/08/2026,date-format=english| and the reversed key order have the
% same interpretation.
%
% |\keys_set:nn| parses and executes the comma-separated family input.
% |\bool_if:NT| executes each commit only when its staging flag is true.
% |\exp_args:NV| passes the stored value, rather than the token-list variable
% itself, to the relevant setter. The function is protected because key
% processing and assignment are not expandable operations. |\ExplSyntaxOff|
% restores the normal catcodes before the public LaTeX2e interface is declared.
%    \begin{macrocode}
\cs_new_protected:Npn \kanban@processsetup #1
  {
    \bool_set_false:N \l__kanban_pending_startdate_bool
    \bool_set_false:N \l__kanban_pending_reference_date_bool
    \keys_set:nn {kanban/setup}{#1}
    \bool_if:NT \l__kanban_pending_startdate_bool
      {\exp_args:NV \kanban@setstartdate
        \l__kanban_pending_startdate_tl}
    \bool_if:NT \l__kanban_pending_reference_date_bool
      {\exp_args:NV \kanban@setreferencedate
        \l__kanban_pending_reference_date_tl}
  }
\ExplSyntaxOff
%    \end{macrocode}
% \begin{macro}{\kanbansetup}
% \begin{macro}{\kbset}
% The canonical setup command exposes the transaction above with one mandatory
% key-list argument. The short alias follows the same collision-tolerant policy
% as the palette alias. This keeps the public namespace convenient without
% taking ownership of an existing command defined by the document class or
% another package.
%
% |\NewDocumentCommand| with signature |m| reads one mandatory argument and
% delegates it unchanged. The LaTeX2e test |\@ifundefined| chooses between
% installing the alias with |\let| and preserving the occupied name with a
% warning. Equivalence by |\let| ensures that |\kbset| and |\kanbansetup| have
% exactly the same argument parser and behavior.
%    \begin{macrocode}
\NewDocumentCommand\kanbansetup{m}{\kanban@processsetup{#1}}
\@ifundefined{kbset}
  {\let\kbset\kanbansetup}
  {\PackageWarning{kanban}
    {\string\kbset\space is already defined;
      \MessageBreak use \string\kanbansetup\space instead}}
%    \end{macrocode}
%
% \paragraph{Collection and application of package options.}
% LaTeX2e presents package options one at a time. The package collects them
% into one token list and applies that list through the same deferred setup
% transaction only after option processing has finished. Package options
% therefore share validation, aliases, key-order independence, and diagnostics
% with |\kanbansetup|.
%
% Formats dated 2021-06-01 or later parse raw option lists with the
% brace-aware key--value processor, so the catch-all |\DeclareOption*|
% collector receives every option intact through |\CurrentOption|, including
% values with embedded commas such as |columns={A,B}|. The older 2020-10-01
% format splits raw option lists at every comma and mishandles brace groups in
% its option bookkeeping, which derails loading. On such formats the package
% routes its options through \textsf{kvoptions} instead: every documented
% setup key is declared as a \textsf{keyval} key whose handler re-wraps the
% brace-stripped value before appending it to the same accumulator, and
% unrecognized options fall through to the same catch-all collector.
%
% |\tl_new:N| allocates the global accumulator. |\tl_if_empty:NF| inserts a
% comma only after the first option, and |\tl_gput_right:Nn| appends globally.
% The expansion prefix |\exp_args:NNo| expands |\CurrentOption| once while
% retaining the destination variable. Application uses |\exp_args:NV| to pass
% the complete accumulated value to |\kanban@processsetup|. The final call to
% |\kanban@applypackageoptions| commits the resulting setup list.
% Unknown package options consequently produce the same warning as unknown
% setup keys on both processing paths.
%    \begin{macrocode}
\ExplSyntaxOn
\tl_new:N \g__kanban_package_options_tl
\cs_new_protected:Npn \kanban@collectpackageoption
  {
    \tl_if_empty:NF \g__kanban_package_options_tl
      {\tl_gput_right:Nn \g__kanban_package_options_tl {,}}
    \exp_args:NNo \tl_gput_right:Nn
      \g__kanban_package_options_tl {\CurrentOption}
  }
\cs_new_protected:Npn \kanban@collectkeyvalue #1#2
  {
    \tl_if_empty:NF \g__kanban_package_options_tl
      {\tl_gput_right:Nn \g__kanban_package_options_tl {,}}
    \tl_gput_right:Nn \g__kanban_package_options_tl {#1={#2}}
  }
\cs_new_protected:Npn \kanban@applypackageoptions
  {
    \tl_if_empty:NF \g__kanban_package_options_tl
      {\exp_args:NV \kanban@processsetup
        \g__kanban_package_options_tl}
  }
\ExplSyntaxOff
\IfFormatAtLeastTF{2021-06-01}
  {%
    \DeclareOption*{\kanban@collectpackageoption}%
    \ProcessOptions\relax
  }
  {%
    \RequirePackage{kvoptions}%
    \SetupKeyvalOptions{family=kanban,prefix=kanban@pkgopt@}%
    \DeclareDefaultOption{\kanban@collectpackageoption}%
    \define@key{kanban}{columns}%
      {\kanban@collectkeyvalue{columns}{#1}}%
    \define@key{kanban}{palette}%
      {\kanban@collectkeyvalue{palette}{#1}}%
    \define@key{kanban}{date-format}%
      {\kanban@collectkeyvalue{date-format}{#1}}%
    \define@key{kanban}{dateformat}%
      {\kanban@collectkeyvalue{dateformat}{#1}}%
    \define@key{kanban}{startdate}%
      {\kanban@collectkeyvalue{startdate}{#1}}%
    \define@key{kanban}{reference-date}%
      {\kanban@collectkeyvalue{reference-date}{#1}}%
    \define@key{kanban}{today}%
      {\kanban@collectkeyvalue{today}{#1}}%
    \define@key{kanban}{mode}%
      {\kanban@collectkeyvalue{mode}{#1}}%
    \define@key{kanban}{columnsep}%
      {\kanban@collectkeyvalue{columnsep}{#1}}%
    \define@key{kanban}{labels}%
      {\kanban@collectkeyvalue{labels}{#1}}%
    \ProcessKeyvalOptions*\relax
  }
\kanban@applypackageoptions
%    \end{macrocode}
% \end{macro}
% \end{macro}
%
%
%    \begin{macrocode}
%</package>
%    \end{macrocode}
